\documentclass[11pt]{article}
\usepackage[a4paper,margin=25mm]{geometry}
\usepackage{amsmath,amssymb,amsthm,booktabs,hyperref}
\hypersetup{colorlinks=true,linkcolor=blue,urlcolor=blue,citecolor=blue}
\newtheorem{theorem}{Theorem}
\newtheorem{lemma}[theorem]{Lemma}
\newcommand{\invGamma}{\operatorname{inv}\Gamma_0}
\newcommand{\HH}{\mathbb H}
\title{Convergence of the inverse-gamma Newton series\\
\large MathDB 357309: proof with an explicit remainder bound}
\author{Principia Math In Progress}
\date{17 September 2026}
\begin{document}
\maketitle

\section*{Task 1: the full conjecture}
Let $\alpha$ be the unique positive zero of the digamma function
$\psi=\Gamma'/\Gamma$, let $\beta=\Gamma(\alpha)$, and let
$\invGamma$ denote the inverse of $\Gamma$ on $[\alpha,\infty)$.
Thus $0<\beta<1$. For $n\geq1$, set
\begin{equation}\label{coeff}
 c_n=\sum_{k=0}^n
 \frac{k}{
 \displaystyle
 \left(\prod_{i=1}^k\log\frac{(k+1)!}{i!}\right)
 \left(\prod_{i=1}^{n-k}\log\frac{(k+1)!}{(k+1+i)!}\right)}.
\end{equation}
Empty products equal $1$.

\begin{theorem}\label{main}
For every real $x>\beta$,
\begin{equation}\label{series}
 \boxed{\displaystyle
 \invGamma(x)=2+\sum_{n=1}^\infty c_n
                    \prod_{i=1}^n\log\frac{x}{i!}.}
\end{equation}
The series converges absolutely and locally uniformly on
$(\beta,\infty)$, with the finite inner sums in \eqref{coeff}
performed before the outer summation.

More explicitly, put $t=\log x$, $a=\log\beta<0$,
$u_0=-a$, $b_j=\log((j+1)!)$, and
\[
 P_N(t)=2+\sum_{n=1}^Nc_n\prod_{j=0}^{n-1}(t-b_j).
\]
For every $N\geq1$,
\begin{equation}\label{bound}
 \left|\invGamma(e^t)-P_N(t)\right|
 \leq
 \frac{|t(t-b_1)|}{b_1(t+u_0)}
       \prod_{j=2}^N\frac{|t-b_j|}{u_0+b_j}.
\end{equation}
The product in this bound is empty when $N=1$.
\end{theorem}

This proves exactly \cite[Conjecture 3.1]{Ulgenes}. Its displayed
sum starts at $n=0$; that term is zero, so \eqref{series} has the
same coefficients and summation convention.

\subsection*{The analytic input}
Define, for real $t>a$,
\[
 F(t)=\invGamma(e^t).
\]
We first justify an integral representation of this particular
inverse of $\log\Gamma$.

We use the classical conformal mapping theorem for $\log\Gamma$
\cite[\S2, equation (2)]{PedersenInverse}, also proved in
\cite[Proposition 6.4]{PedersenSurvey}. For the branch of $\log\Gamma$
real on the positive axis, the upper half-plane $\HH$ is mapped
conformally onto
\[
 \mathcal V=\mathbb C\setminus
       \bigcup_{j=0}^\infty
       \bigl([\log|\Gamma(\alpha_j)|,\infty)-ij\pi\bigr).
\]
Here $\alpha_0=\alpha$, and for $j\geq1$ the point
$\alpha_j\in(-j,-j+1)$ is the unique zero of $\psi$ in that interval.
Every removed ray has nonpositive imaginary part, so
$\HH\subset\mathcal V$. It follows that the inverse of $\log\Gamma$
restricted to $\HH$ maps $\HH$ into $\HH$: it is a Pick function.

For real $t>a$, this inverse has the real boundary value
$F(t)>\alpha$. Indeed, $\log\Gamma$ is strictly increasing on
$(\alpha,\infty)$ and its derivative there is positive; its local
inverse therefore takes nearby points above the real axis to points
above the real axis. Uniqueness of the conformal inverse identifies
the branches. The local inverse theorem and Schwarz reflection
extend $F$ analytically across all of $(a,\infty)$, with real values.

The Herglotz representation and its support characterization
\cite[p.~90]{PedersenSurvey} consequently give
\begin{equation}\label{repr}
 F(t)=A t+B+
       \int_{u_0}^\infty
       \left(\frac{u}{1+u^2}-\frac1{t+u}\right)\,d\rho(u),
 \quad t>a,
\end{equation}
where $A\geq0$, $B\in\mathbb R$, $\rho$ is a positive measure,
and $\int_{u_0}^\infty(1+u^2)^{-1}\,d\rho(u)<\infty$.
The support follows from real analytic continuation across
$(a,\infty)$, followed by the change of variable $u=-s$.
Both the conformal mapping theorem and the Herglotz theorem are
established inputs, explicitly credited here.

Since $b_0=0$, $b_1=\log2$, $F(0)=2$, and $F(b_1)=3$,
subtracting \eqref{repr} at these two points yields
\begin{equation}\label{moment}
 \frac1{b_1}
   =A+\int_{u_0}^\infty\frac{d\rho(u)}{u(u+b_1)}.
\end{equation}
In particular the integral on the right is at most $1/b_1$.

\subsection*{Identifying all coefficients}
Let $I_NF$ be the polynomial of degree at most $N$ interpolating
$F$ at $b_0,\ldots,b_N$. These are distinct nodes, and
\[
 F(b_j)=j+2.
\]
The Newton coefficient of order $n\geq1$ is the divided difference
\[
 F[b_0,\ldots,b_n]
   =\sum_{k=0}^n\frac{k+2}{\prod_{\substack{0\leq j\leq n\\j\ne k}}(b_k-b_j)}
   =\sum_{k=0}^n\frac{k}{\prod_{\substack{0\leq j\leq n\\j\ne k}}(b_k-b_j)}.
\]
The constant $2$ contributes zero because a constant has zero
divided differences of positive order; equivalently, compare the
leading coefficient in its Lagrange interpolation formula.
For $j<k$ and $j>k$, respectively, the denominator factors are
\[
 b_k-b_j=\log\frac{(k+1)!}{(j+1)!}.
\]
Reindexing the first group by $i=j+1$ and the second by $i=j-k$
gives exactly \eqref{coeff}. Thus $I_NF=P_N$, not merely a formal
series with matching first coefficients.

\subsection*{The exact remainder}
For fixed $u\geq u_0$ let $h_u(t)=-1/(t+u)$. Its interpolation
remainder has the rational identity
\begin{equation}\label{kernel}
 h_u(t)-I_Nh_u(t)
   =\frac{(-1)^N\prod_{j=0}^N(t-b_j)}
          {(t+u)\prod_{j=0}^N(u+b_j)}.
\end{equation}
To verify it, multiply the left side by $t+u$.
The result is a polynomial of degree at most $N+1$ vanishing at
every $b_j$ and taking the value $-1$ at $t=-u$.
These conditions give precisely the right-hand numerator and sign.

For $N\geq1$, interpolation reproduces the affine part of
\eqref{repr}, as well as the constant term in each regularized
integrand. A finite Lagrange sum can be interchanged with the
convergent regularized integrals. It follows that
\begin{equation}\label{remainder}
 F(t)-P_N(t)=(-1)^N\prod_{j=0}^N(t-b_j)
       \int_{u_0}^\infty
       \frac{d\rho(u)}{(t+u)\prod_{j=0}^N(u+b_j)}.
\end{equation}
No separately divergent constant integral has been split off.

For $t>a$ all the denominator factors in \eqref{remainder}
are positive. Extract the factors with $j=0,1$ and use
$u\geq u_0$, $t+u\geq t+u_0>0$, and \eqref{moment}.
This proves \eqref{bound}:
\[
 |F(t)-P_N(t)|
 \leq
 \frac{|t(t-b_1)|}{t+u_0}
 \left(\prod_{j=2}^N\frac{|t-b_j|}{u_0+b_j}\right)
 \int_{u_0}^\infty\frac{d\rho(u)}{u(u+b_1)}.
\]
It also proves the one-sided remainder rule
\begin{equation}\label{signerror}
 (-1)^N\prod_{j=0}^N(t-b_j)\,
       (F(t)-P_N(t))\geq0.
\end{equation}

\subsection*{Convergence throughout the stated interval}
Fix $t>a$. Once $b_j>t$, the product factor in \eqref{bound} is
\[
 \frac{|t-b_j|}{u_0+b_j}
   =1-\frac{t+u_0}{u_0+b_j}\quad\hbox{in }(0,1).
\]
Furthermore,
\begin{equation}\label{divergence}
 \sum_{j=2}^\infty\frac1{u_0+b_j}=\infty,
\end{equation}
because $b_j=\log((j+1)!)\leq(j+1)\log(j+1)$ and
$\sum_{j\geq2}1/(j\log j)$ diverges. The inequality
$1-v\leq e^{-v}$ now shows that the product in \eqref{bound}
tends to zero. If $t$ is a node, interpolation is already exact
once that node is included. This proves $P_N(t)\to F(t)$ for
every $t>a$, including negative $t$.

The convergence is uniform on a compact interval $K\subset(a,\infty)$.
Choose $J$ such that $b_j>\max K$ for $j\geq J$.
The finitely many initial factors and the prefactor in
\eqref{bound} are uniformly bounded on $K$. In the remaining
factors, replace $t+u_0$ by $\min K+u_0>0$ and apply
\eqref{divergence}. This proves local uniform convergence.

For $n\geq2$, taking divided differences in \eqref{repr} gives
\begin{equation}\label{signcoef}
 c_n=(-1)^{n-1}
       \int_{u_0}^\infty\frac{d\rho(u)}
                              {\prod_{j=0}^n(u+b_j)}.
\end{equation}
Thus the coefficient signs alternate weakly. For fixed real $t$,
all sufficiently late Newton terms have the same sign, since
$t-b_j<0$ for all sufficiently large $j$. The already established
convergence is therefore absolute. This is also locally uniform
absolute convergence: after an index independent of $t\in K$,
the absolute tail at each $t$ equals $|F(t)-P_N(t)|$, which
tends uniformly to zero. Finally, the substitution $t=\log x$
proves Theorem \ref{main} in full.

\subsection*{Verification and scope}
The exact verifier checks the rational-kernel base case,
divided-difference induction, remainder update, and regularization
cancellation, as well as independent Lagrange constructions at
six finite node sets. There are sixteen exact symbolic identities.
Independent recursive and printed-product coefficients agree in
24 cases, and 128 computed coefficients have the predicted signs.

Eleven real test points, including four factorial nodes, are checked
at degrees $8,32,128$ against a separately computed principal
inverse and the bound \eqref{bound}, using 220-digit arithmetic
with a 160-digit precision comparison. All 33 bound diagnostics pass.
These checks do not prove convergence; the all-$x$ argument is above.
Convergence can be slow: at $x=0.89$, the degree-128 error is about
$0.182$. The theorem asserts convergence, not rapid numerical accuracy.

The complete published source and its latest arXiv version both
retain Conjecture 3.1 on page 9. Dated bounded status and literature
checks accompany this note. The authoring agent performed the audit;
no independent peer review is claimed.

\newpage
\section*{Tasks 2--5: proposed follow-up work}
These are four proposed tasks, not additional solved problems.
\begin{enumerate}
 \setcounter{enumi}{1}
 \item \textbf{Task 2.} Determine the leading asymptotic of
 $F(t)-P_N(t)$ for each fixed nonnode $t>a$, using the behavior
 of the representing measure near $u_0$; compare it with
 the product bound \eqref{bound}.
 \item \textbf{Task 3.} Determine convergence on the boundary
 $\Re t=a$ for the corresponding complex Newton series, including
 its value at the real endpoint $t=a$.
 \item \textbf{Task 4.} Construct a numerically stable accelerated
 evaluation of \eqref{series}, with a certified error tolerance
 on any prescribed compact subinterval of $(\beta,\infty)$.
 \item \textbf{Task 5.} Determine the limit, if it exists, of the
 Newton interpolation series for $\invGamma(x)$ at the unlogged
 factorial nodes $(j+1)!$, and compare that limit with $\invGamma$.
\end{enumerate}

\begin{thebibliography}{9}
\bibitem{Ulgenes} D. P. H. Ulgenes,
\emph{Series and Product Representations of Gamma, Pseudogamma and
Inverse Gamma Functions}, Maple Transactions \textbf{5}(1),
Article 16090 (2025), 10 pp.,
\href{https://doi.org/10.5206/mt.v5i1.16090}{doi:10.5206/mt.v5i1.16090}.
Also \href{https://arxiv.org/abs/2301.09699v4}{arXiv:2301.09699v4}.
\bibitem{PedersenInverse} H. L. Pedersen,
\emph{Inverses of gamma functions},
\href{https://arxiv.org/abs/1309.2167}{arXiv:1309.2167}, 2013,
\S2, equation (2).
Published in Constructive Approximation \textbf{41} (2015), 251--267.
\bibitem{PedersenSurvey} H. L. Pedersen,
\emph{Pick Functions as a Tool in the Theory of Special Functions},
RIMS K\^oky\^uroku \textbf{1893} (2014), 89--101,
\href{https://www.kurims.kyoto-u.ac.jp/~kyodo/kokyuroku/contents/pdf/1893-12.pdf}{complete published survey}.
See p.~90 for Herglotz representation and support, and
pp.~99--100 for the conformal mapping theorem.
\end{thebibliography}
\end{document}
